%%%%%%%%%%%%%%%%%%%%%%%%%%%%%%%%%%%%%%%%%%%%%%%%%%%%%%%%%%%
% EPFL report package, main thesis file
% Goal: provide formatting for theses and project reports
% Author: Mathias Payer <mathias.payer@epfl.ch>
%
% To avoid any implication, this template is released into the
% public domain / CC0, whatever is most convenient for the author
% using this template.
%
%%%%%%%%%%%%%%%%%%%%%%%%%%%%%%%%%%%%%%%%%%%%%%%%%%%%%%%%%%%
\documentclass[a4paper,11pt,oneside]{report}
% Options: MScThesis, BScThesis, MScProject, BScProject
\usepackage[MScProject,lablogo]{EPFLreport}
\usepackage{xspace}
\usepackage{xcolor}
\usepackage{listofitems}
\usepackage{subcaption}
\usepackage{booktabs}
\usepackage{longtable}

\newif\ifreview
\reviewtrue
% To remove the reviews seet \reviewfalse
%\reviewfalse

\definecolor{amber}{rgb}{1.0, 0.75, 0.0}
\newcounter{ReviewerID}
\readlist*\annotationcolors{blue, red, orange, green, purple, amber, brown, olive}
\newcommand{\newreviewer}[2]{%
  \ifnum \theReviewerID=\annotationcolorslen
  \setcounter{ReviewerID}{0} % cycle through colors
  \fi
  \stepcounter{ReviewerID}%
  \expandafter\edef\csname bootstrap#1\endcsname{%
    \expandafter\def\csname #1\endcsname####1{%
      \ifreview%
	  {\noexpand\color{\annotationcolors[\theReviewerID]} {\noexpand\bf{\noexpand\fbox{#2}} {\noexpand\it ####1} }}
	  \else%
	      {}% disable annotations
	      \fi%
    }%
  }%
  \csname bootstrap#1\endcsname%
}

\newreviewer{philipp}{philipp}
\newreviewer{elias}{elias}

\title{House of Scudo}
\author{Elias Valentin Boschung}
\supervisor{Mao Philipp Yuxiang}
\adviser{Prof. Dr. sc. ETH Mathias Payer}
%\coadviser{Second Adviser}
%\expert{The External Reviewer}

\newcommand{\sysname}{House of Scudo\xspace}

\begin{document}
\maketitle{}
\dedication{
  \begin{raggedleft}
    If debugging is the process of removing software bugs,\\
    then programming must be the process of putting them in.\\
    --- Edsger Dijkstra\\
  \end{raggedleft}
  \vspace{4cm}
  \begin{center}
    Dedicated to my study companion since Covid-19, my plush bear.
  \end{center}
}
\makededication{}


\acknowledgments{I thank first and foremost my family for their unconditional support
  throughout all of my studies until this point. I also thank my supervisor Philipp for
  all the tips and help in writing this project.  Thanks to my friends, those who inspired
  and encouraged me on the path of cybersecurity, as well as those who helped me relax
  when I was stressed.  Thanks also to the HexHive lab for the opportunity to do this
  Master Semester Project, and also for providing the template for this report, which is
  available at \url{https://github.com/hexhive/thesis_template}.}  \makeacks{}

\begin{abstract}
  Heap bugs are a very frequent issue in C programming, due to the quite high
  complexity of memory management for programmers. Attackers can exploit such heap bugs
  like overflows, double free's and similar to abuse the underlying mechanics of the heap
  allocator and get even bigger control over heap memory such as arbitrary writes. For the
  common libc allocator there are a lot of known techniques to achieve that, a few of
  which are known as heap exploitation houses. Scudo is an allocator used on
  Android 11+ that is supposed to be specifically hardened to prevent or at least
  complicate such attacks.~\cite{llvmScudo} However, there exists no documentation of any
  tests of the resilience of Scudo to a dedicated attacker.

  In this project we looked for exploits using these heap vulnerabilities and
  targeting the mechanics of Scudo specifically. We discovered two exploits that use
  different heap vulnerabilities and mechanics of the Scudo allocator to trick it into
  returning a chunk at an address chosen by the attacker. These two exploits therefore
  enable an arbitrary write primitive similar to the primitives provided by some libc
  exploitation houses.
  
  In addition to the attacks, we propose a mitigation that prevents either of the
  discovered exploits by targeting a mechanic of the Scudo allocator that is used in both
  of our exploits and which further hardens Scudo. We show that the impact of our
  mitigation on the allocator's performance should be negligible in most common scenarios.
\end{abstract}

\maketoc{}

%%%%%%%%%%%%%%%%%%%%%%
\chapter{Introduction}
%%%%%%%%%%%%%%%%%%%%%%

Heap vulnerabilities are a frequent and dangerous attack surface and have been for a long
time, as even Microsoft already stated back in 2013 that they were the most common type of
vulnerability that they addressed through security updates.~\cite{microsoftHeapReport}
These vulnerabilities can be used by an attacker to get remote code execution, essentially allowing
the attacker to take control of the device.
Attackers often leverage the underlying mechanisms of the heap allocator itself.
For the widely used glibc allocator there exist
generic techniques that, given some preconditions, give a pattern to escalate a heap
vulnerability and to get remote code execution. 
like the unsafe unlink exploit or
the list of libc heap exploitation houses. Such heap vulnerabilities are also relevant in
the Android environment, as can be seen by the example of a remote code execution found in
WhatsApp back in 2019, which was caused by a double free.~\cite{whatsappRCE}

To harden applications on Android against such attacks, it does not use the popular glibc
allocator, but since Android 11 uses a new heap allocator called Scudo.~\cite{llvmScudo}
Scudo was developed specifically with security in mind, to reduce the possibilities of an
attacker to leverage the allocators mechanisms to get an arbitrary write primitive. Scudo
uses a checksum to secure critical metadata stored in-line and it also tries to isolate
some of the metadata, specifically all the free list related metadata, to its own region,
to reduce the risk of possible corruption with a simple overflow. Since it is designed
completely differently, there is currently no existing study on an attack that uses the
known techniques for the glibc allocator to exploit Scudo.

Therefore, in this project we try to find techniques to exploit mechanisms from the heap
allocator itself, similar to the way such techniques exist for the glibc heap
allocator. We start by defeating the checksum with a single leak of the metadata of a
chunk, including the corresponding checksum. Then we show two exploits, that build on
defeating the checksum to forge different headers and tricking Scudo to process them
without any errors, which will allow us to gain an arbitrary write. Additionally to these
two exploits, we also build some additional tools to facilitate further research and
exploiting of Scudo, such as a python library to break the checksum and some templates for
different exploit approaches to Scudo.

Finally, we detect that the two attacks we found both rely on a common mechanism, being how
Scudo handles larger chunks compared to smaller ones. In order to protect against these
attacks, we propose a mitigation that adds some additional checks to address that specific
exploitation of the mechanism, and thus prevents these attacks. We evaluate our mitigation
and find that the overhead on benchmarks is negligible. We therefore open a pull request
to llvm, the official development team behind the Scudo allocator.

%%%%%%%%%%%%%%%%%%%%%%%%%%%%%%%%%
\chapter{Threat model}
%%%%%%%%%%%%%%%%%%%%%%%%%%%%%%%%%

In our threat model, the attacker has access to some basic heap vulnerabilities, such
as overflows, double free's and arbitrary free's. Additionally, if address space layout
randomization is used (ASLR), the attacker also has access to an ASLR leak through a
channel, allowing him to know where some memory regions are loaded.
These heap vulnerabilities can give the
attacker some control over a restrained space of heap memory. The goal of the attacker
however is to get an arbitrary write primitive, them to get full control of the 
program. We don't consider any vulnerabilities in application specific data, so the
attacker from our threat model should only exploit the allocator's mechanisms and not rely
on any application data, such as data or code pointers stored on the heap.

When using the glibc allocator, or some other unhardened allocators, these basic heap
vulnerabilities can directly lead to an arbitrary write primitive, by using pointers
stored in the controllable heap memory. In glibc this can for example be achieved by using
a tcache forward pointer corruption, which exploits the fact that the cache linked list is
stored inside the free chunks themselves. Since these free chunks may be controlled by a
use after free or an overflow for example, the cache can be corrupted and an arbitrary
address inserted. An example of such a corruption can be found on~\cite{tcachePoisoning}.

However, Scudo is built with such security issues in mind, therefore the use of pointers in
the same memory regions as the chunks with data is avoided as much as possible. Therefore
the attacker needs to trick some other, more complicated allocator mechanisms in order to
gain the full arbitrary read or write primitive.

We do not consider any attacks on the quarantine, since it is not used in the
configuration of Scudo used on Android devices. Even on the llvm website, the development
team behind Scudo, they mention that quarantine may have a quite high impact on
performance and is therefore disable by default.

%%%%%%%%%%%%%%%%%%%%%%%%%%%%%%%%%
\chapter{Scudo Security Measures}
%%%%%%%%%%%%%%%%%%%%%%%%%%%%%%%%%

\philipp{right now this section isn't very focused on security measures}

\philipp{you should explicitly introduce scudos two main security mechanism and 
then have one section for each (after the basics section)}

\philipp{one section about protected inline metadata: only very little inline 
metadata. free list in segregated region/bss of scudo, 
headers protected with checksum, secondary chunk header metadata, 
stored in front of mmapped regions}

\philipp{another section about segregated memory regions and random allocations 
inside those regions (shorter than previous one, since we don't care too 
much about this for our exploits)}

\philipp{we will introduce all the other mechanisms when they become relevant 
in the exploits/mitigations (such as the free list for the chunks)}

\section{Basics}

\philipp{don't talk about performance here, rest is fine}
The Scudo hardened allocator divides the heap allocations it is asked to do into two big
types of allocations, based on the size of the chunk to be allocated.  The smaller chunks
are handled by the primary allocator inside Scudo, which is designed to optimize
performance as much as possible by adapting to finer-grained size differences between
these smaller chunks. All the big chunks are handled by the secondary allocator, which is
less optimized, since the most frequent chunks are the small ones and thus the primary
allocator has a bigger impact on performance than the secondary allocator.

\section{Combined Chunk Header}

\begin{figure}[h]
  \centering
  \begin{tabular}{l r}
    \toprule
    ClassId & 8 bits \\
    State & 2 bits \\
    OriginOrWasZeroed & 2 bits \\
    SizeOrUnusedBytes & 20 bits \\
    Offset & 16 bits \\
    Checksum & 16 bits \\
    \bottomrule
  \end{tabular}
  \caption{Layout of the Scudo chunk header}\label{fig:ScudoChunkHeader}
\end{figure}

The chunks allocated by both the primary and the secondary allocator store some metadata
in a combined header. This header is then stored just in front of the actual content of
the chunk, such that it can easily be checked whenever the chunk gets accessed. An
overview of the contents of this header can be seen in \autoref{fig:ScudoChunkHeader}. The
ClassId identifies the class the chunk belongs to in case it was allocated by the primary
allocator, with a ClassId of 0 meaning that the chunk was allocated by the secondary
allocator. The header also contains information about the state of the chunk, which can be
allocated, quarantined or available, as well as the origin of the allocation, e.g., malloc
or new, which can be used to detect errors when the type of deallocation does not match
the type of the allocation. Furthermore, the header includes the size for the primary
chunks or the number of unused bytes for the secondary chunks and an offset, which is the
distance from the beginning of the returned chunk to the beginning of the actual backend
allocation.  Finally, the header contains a checksum, which is generated using a cookie (a
random number generated during the initialization of the allocator), the heap address of
the chunk, and the actual content of the header. This checksum is used for guarding
against programming errors as well as attackers, and it is checked each time the header is
loaded.

\section{Secondary Chunk Header}

\begin{figure}[h]
  \centering
  \begin{tabular}{l r}
    \toprule
    Prev & 64 bits \\
    Next & 64 bits \\
    CommitBase & 64 bits \\
    CommitSize & 64 bits \\
    MemMap & 128 bits \\
    \bottomrule
  \end{tabular}
  \caption{Secondary chunk header on linux}\label{fig:ScudoSecondaryChunkHeader}
\end{figure}

The secondary header stores some additional information in a header which is prepended to
the normal chunk header. First it stores pointers to the next and previous large chunk
headers, which forms the doubly linked list of allocated secondary blocks. These pointers
are used in the Safe Unlink Secondary exploit which is explained in a later section. Then
it stores the CommitBase and CommitSize, which are the address of the chunk itself and its
size. Lastly it stores some mapping information which can depend on the underlying
operating system, but in the case of a Linux OS it simply corresponds to the MapBase and
CommitBase. These are similar to CommitBase and CommitSize, but they may be slightly
different due to alignment requirements. All these values are, contrary to the combined
chunk header, not protected by a checksum, so they could be changed without trouble if we
have write access to them.

\section{Secondary Cache}

The secondary allocator has its own cache, in which it simply stores the free'd secondary
chunks in a simple list, and before allocating a new chunk it searches the list for a
chunk whose size is at most a certain number of pages larger than the requested size. The
values stored in this cache list, are based on the values of the secondary chunk header at
the moment the chunk is free'd. Interestingly, no integrity checks are done on the values
saved in the cache, which is used in the CommitBase exploit which is explained in a later
section.


%%%%%%%%%%%%%%%%%%%%%%%%%%%%%
\chapter{Breaking the Cookie}
%%%%%%%%%%%%%%%%%%%%%%%%%%%%%

To check the integrity of chunk headers, Scudo uses a simple checksum with a random 32-bit
cookie which is used as seed to the hash algorithm. It may use one of two hashing
algorithms, either CRC32 if hardware support is available for it, or the BSD checksum
algorithm. The following pseudo code demonstrates how the checksum is calculated.
\begin{verbatim}
// Inputs
uint32 cookie = randomInt()
uint64 address
uint64 header

// Output
uint16 checksum

// Logic
If USES_HARDWARE_CRC32:
  uint32 intermediate = CRC32(cookie, address)
  intermediate = CRC32(intermediate, header)
  checksum = (intermediate & (intermediate >> 16)) & 0xffff
Else:
  uint16 seed = cookie & 0xffff
  uint16 intermediate = BSD(seed, address)
  intermediate = BSD(intermediate, header)
  checksum = intermediate
\end{verbatim}

This checksum protects against blind heap overflows, since in a blind heap overflow the
attacker has no way to forge the correct checksum. Therefore, if the attacker tries to
modify any values in the header, Scudo will try to verify the checksum of the header
before taking any action with the chunk, and will crash the program since the checksum
will not match the data in the header.

However, these checksum/hash algorithms are not designed to be cryptographically
secure. While the cookie is 32 bits long, the checksum is only 16 bits, and it is
therefore possible to find a cookie collision and produce the same checksums the original
cookie would have produced. Indeed it is enough to get a single heap leak of a complete
chunk header together with its address in order to recover the cookie 
and forge valid chunk headers.
Recovering the cookie can be done by simply bruteforcing the
cookie value until we find the same checksum that we leaked. Such a bruteforce of the
cookie was already used by the un1fuzz hobby research group, and another researcher used a
z3 solver to achieve the same thing.~\cite{un1fuzz}~\cite{z3SolveScudoCookie}

Using a C program bruteforcing the cookie takes around 1–2 seconds, therefore we
implemented a C library that takes care of the checksum calculation and brute-forcing of
the cookie given the necessary data. To facilitate integration with the pwntools python
tooling often used for exploit development, we also provide python bindings for this
library.

%%%%%%%%%%%%%%%%%%
\chapter{Exploits}
%%%%%%%%%%%%%%%%%%

\section{Change Class ID}

This exploit uses the breaking of the checksum to modify some of the inline header
metadata, to trick the allocator into processing a chunk differently. It is used in the
two other exploits, which then take further advantage of the processing of the chunk that
is different from the intended way.

\subsection{Summary}

Requirements:
\begin{itemize}
\item Header leak with corresponding address
\item Free where we control the \emph{0x10} bytes in front of the address
\end{itemize}

Result: Move a primary chunk to another size class, making potentially larger chunks in a
small chunk region and multiple chunks overlap.

\subsection{Requirements}

The change class ID exploit needs some flaws in the target binary to be applicable.  First
we need to somehow leak a header value along with its address, in order to brute-force the
cookie and forge new header checksums like explained earlier. Then the attacker also needs
to have a free operation occurring, at an address in front of which we can write some
bytes, specifically \emph{0x10} bytes in front of the free target. Finally, there needs to
be some primary chunk allocations after the free, which are located in the same class ID
as our target we want to move the chunk to.

\subsection{Explanation}

For the retrieval of the cookie refer to the previous section for a thorough explanation
on how it works.

Once the attacker has the cookie, he needs to find a vulnerable free operation, where he
can write data to the \emph{0x10} bytes in front of the free'd address. 
This can be achieved in different ways, such as an overflow while writing to a chunk of a
small size, causing the attacker to be able to write over the header of one of the next
chunks. Another possible way to achieve the primitive can be controlling a pointer that is
freed, and setting that pointer to a memory region that was set up to look like a fake
chunk, thus freeing the completely fake chunk.

Though the exact offset from the first chunk which contains the overflow
to the target header is random due to Scudo's defense mechanisms, the attacker can still
manage to overwrite it, since the offset is a multiple of a fixed size that depends on the
size class that is attacked. The target header might still be allocated before the chunk
which the attacker can overflow, but with some random luck and if the attacker can execute
the program as often as he wants, at some point the target header should be allocated at
one of these multiples after the overflow.

The attacker then forges a new valid header with its corresponding checksum thanks to
the brute-forced cookie. 
In this case the attacker probably wants to get a large chunk in
the small chunk region, thus he will replace the header with one specifying a larger class
ID instead of the original one.

When the chunk with the forged header gets free'd, it will be put into the cache of the
larger class ID, and when a chunk from that same class ID gets allocated again, Scudo will
allocate a larger sized chunk in the small chunk region.

\begin{figure}[h!]
  \centering
  \includegraphics[width=\linewidth]{figures/ChangeClassID.png}
  \caption{Illustration of the modification to chunk header}\label{fig:ChangeClassID}
\end{figure}

\philipp{some diagram/graphic etc to show how the exploit works, what happens 
when it succeeds etc would be very helpful for the reader}

\section{Exploit CommitBase}

In this exploit, the attacker modifies the unprotected metadata of a fake secondary chunk,
to put an address of the attacker's choosing into the cache for secondary chunks. This
injected address is then used to directly allocate a new secondary chunk at the address
chosen by the attacker. The exploit relies on the Change Class ID exploit to get a fake
secondary chunk, since real secondary chunks are protected against overflows with guard
pages.

\subsection{Summary}

\begin{itemize}
\item Header leak with corresponding address
\item At least one secondary chunk allocated
\item Free on any address where we control \emph{0x40} bytes in front
\item Allocation of secondary chunk of known (approximate) size
\end{itemize}

Result: We can choose an arbitrary location for the secondary chunk allocated

\subsection{Requirements}

The forge secondary exploit has a certain number of requirements that the target binary
needs to fulfill.

First we need a header leak of any chunk from the heap as well as its address, in
order to brute-force the cookie. Afterward we can forge any header checksum from the
cookie we calculated from this leak.

Second we need a free for which we can control a certain number of bytes before the free,
more specifically we need to control \emph{0x40} bytes in front of the address of the
free. This is the size of the secondary header plus the size of the padded primary header.

Third we need there to be at least one secondary chunk allocated before the free, and we
need there to be another secondary allocation afterward, of which we know the approximate
size. This is the allocation where we will get our arbitrary access.

\subsection{Explanation}

\philipp{figure as well}
The breaking of the cookie works exactly the same as explained previously, and the first
part of this exploit also reuses a previously introduced exploit. Indeed, the first part is
basically a change class ID exploit, except that we don't change the class ID to a
different primary one, but to the special class ID 0, that is used to designate chunks
from the secondary allocator.

The second part exploits the unsecured secondary header and the secondary cache. When
freeing a secondary chunk, the info from its secondary header is put into the cache, and
when allocating a new secondary chunk the cache is checked for a similarly sized
space. Since we can change the secondary header of our fake secondary chunk, we can change
the CommitBase and MapBase to any address where we want to allocate a chunk, and
CommitSize and MapSize to our desired size. When setting the size we need however to pay
attention to the fact that the cache is only checked for similar sized chunks, so the size
we set can only be a certain number of pages larger than the next allocation where we want
our chunk to be allocated. The exact allowed size difference may depend on the specific
configuration of Scudo.

Following that, the next secondary chunk should be located at our chosen address. The
special requirement that there was at least one secondary chunk allocated previously, is
due to the way the free of a secondary chunk works. Indeed, freeing a fake secondary chunk
that was not correctly allocated previously somewhat corrupts the tracking of allocated
secondary chunks. While the list of allocated chunks itself is stored as a linked list
with the chunks itself, and as such the freeing of the fake secondary chunk doesn't touch
that list, it still modifies the number of items recorded for the list. Since this number
is stored as an unsigned int, if there was no previously allocated secondary block, our
free of the fake secondary chunk would cause a negative overflow, and on the next
allocation the program would crash, so before we could get our arbitrary chunk.


\section{Safe Unlink Secondary}

This exploit uses a trick similar to the libc safe unlink exploit, that tricks the linked
list deletion function into inserting a pointer to Scudo's free list into the free list
itself. The attacker achieves this by building a fake linked list, and is rewarded with
the ability to write into the internal Scudo structures, which can be escalated into an
arbitrary write.

\subsection{Summary}

Requirements:
\begin{itemize}
\item Header leak with corresponding address
\item Scudo library base leak
\item Three consecutive frees on the same address
\item Control over the \emph{0x10} bytes in front of the address before the first two
  frees
\item Control over the \emph{0x40} bytes in front of the address before the third free
\end{itemize}

Result: Next allocation will be a chunk in the PerClass structure itself (thread-local
free list), allowing control of future allocations


\subsection{Requirements}

The safe unlink exploit has a certain number of requirements that the target binary needs
to fulfill.

First we need a header leak of any chunk from the heap as well as its address, in
order to brute-force the cookie. Afterward we can forge any header checksum from the
cookie we calculated from this leak.

Second we need a free for which we can control a certain number of bytes before the free,
more specifically we need to control \emph{0x40} bytes in front of the address of the
free. This is the size of the secondary header plus the size of the padded primary header.

Third we want to forge a fake linked list, for which we need a leak of the Scudo library
base, and two more frees of the chunk in the previous step and
being able to control the \emph{0x10} bytes in front of the address to modify the
header.

\subsection{Explanation}

\begin{figure}[h!]
  \centering
  \includegraphics[width=\linewidth]{figures/SafeUnlinkMemoryFrees.png}
  \caption{Illustration of forged linked list and the frees we
    need}\label{fig:SafeUnlinkMemoryFrees}
\end{figure}

As with most Scudo exploits, the first step is to get the cookie, which can be done from
any single header leak, as explained in the dedicated section. This willl be used later,
to craft a fake secondary chunk by changing the class ID saved in the header to 0. This
causes Scudo to look at the forged secondary chunk header and to parse some of the data
contained in it, like the pointers to the previous and next in use secondary chunks. In
this exploit we target the linked list pointers and the logic to remove an element from
the linked list. An overview of what our forged linked list should look like can be seen
in \autoref{fig:SafeUnlinkMemoryFrees}. In that figure are also shown the locations of the
three different frees that are needed and that are explained in more detail below.

Then we need to have two locations in memory that point to the chunk we will tamper with
in the next section. The way we achieve this is by having two frees on that same chunk,
leaving its address twice in the PerClass structure next to each other. The PerClass
structure constitutes the first level, thread specific free list of chunks. It is handled
as a simple array of chunk addresses, and therefore freeing the same chunk twice leaves
two pointers to that chunk right next to each other in the PerClass structure. In order to
not trigger double free detection, we just need to modify the header of the chunk between
the two frees to set its status to Allocated again. Since Scudo completely trusts the
header and doesn't try to detect a double free in any other way, this is sufficient to get
the double free without Scudo throwing any error.

\begin{figure}[h!]
  \centering
  \includegraphics[width=\linewidth]{figures/SafeUnlinkMemory.png}
  \caption{Illustration of different important memory regions}\label{fig:SafeUnlinkMemory}
\end{figure}

Once we have set up the two addresses in the PerClass array, we need to know its address
still. To calculate the address of the PerClass array(s) we need to have a leak of some
Scudo address, like the base address where Scudo is loaded or the address of the malloc
function. We can calculate the offset needed from there based on the Scudo version, and
then advance by \emph{0x100} times the value of the class ID where we free'd the
chunks. Then we just need to guess the number of chunks present in the PerClass structure
of that class ID. We will refer to the addresses that we calculate in this step as the
\emph{perclass\_addr1} and \emph{perclass\_addr2} from here on forward for simplicity. These
two addresses both point to a location in the PerClass array, at which our two addresses
of our chunks are located.

Finally, we need to prepare the header of the final free that will trigger the safe
unlink. For that we need to configure a fake secondary header in front of the primary
header, as well as modifying the latter to set its class ID to 0. We set the prev
(\emph{chunk\_addr - 0x40}) to \emph{perclass\_addr2 - 0x8} and the next
(\emph{chunk\_addr - 0x38}) to \emph{perclass\_addr1}. With the PerClass setup, we
can set both of them to the address where the first of the two addresses is stored.  Next
follow the CommitBase, CommitSize and MapBase, MapSize. We can set them to the values we
want, we can set them to whatever since we don't actually use them. They have a size of
\emph{0x8} bytes each.  Finally, we just need to set the primary header with the class ID
set to 0 and the checksum recalculated.

When the free of that chunk then happens, it is handed to the secondary allocator since we
set the class ID to 0. The secondary allocator tries to remove the chunk from the linked
list of in use chunks, even when our fake chunk never was in it. So it tries to set the
Next$\rightarrow$Prev pointer to Prev and the Prev$\rightarrow$Next pointer to Next, which
leads it to write the addresses of our locations in the PerClass structure
(\emph{perclass\_addr1} and \emph{perclass\_addr2}) to those same addresses, which will
allow us to allocate a chunk in the PerClass structure itself.

%%%%%%%%%%%%%%%%%%%%
\chapter{Mitigation}
%%%%%%%%%%%%%%%%%%%%

Both our exploits depend on the same mechanism of how secondary chunks are handled, by setting
the ClassId of a chunk to 0 before freeing it and therefore tricking Scudo into handling a
small chunk as if it was a big chunk. In the process of freeing the fake big chunk, Scudo
then parses metadata before the chunk header that is controlled by the attacker.
Therefore, in our mitigation we implement an additional system to verify before freeing a
large chunk that it was actually allocated as a large chunk.

To be able to do this verification, we allocate a new special memory region, which can by
default hold 1001 addresses. Whenever a new big chunk is allocated, its address is stored
in the first free space in this special memory region, where a free space is defined by
the address being 0. If the 1000 first addresses in this special memory region are all
already used, a new region of 1001 addresses is allocated, and the address of this new
region is stored as the 1001st address in the first region, creating space for 1000
additional addresses. 

Whenever a big chunk is requested to be free'd, this special memory region is first
searched for its address. If it is found, the space in the special memory region is set to
0 to indicate it is free again, and the rest of the logic to free the big chunk continues
normally. If however the address is not found in the special memory region, the free is
ignored as it is assumed to be either a corrupt chunk or an attacker trying to trick Scudo
into freeing a fake forged chunk.

Two options were added in the Allocator Config to configure the details of the mitigation
and to completely disable the mitigation if it is deemed too expensive in some
configurations.

\section{Evaluation}

\subsection{Effectiveness}

After implementing the mitigation, we first checked to see if our exploits were indeed
fixed by our mitigation. Since our fake chunks we are trying to free to trick the
allocator mechanisms were in both cases originally small chunks, the mechanisms are never
even triggered since our mitigation detects that the fake chunks are not supposed to be
large chunks. Therefore Scudo does abort the processing of the chunk before it gets to
parsing the attacker controlled metadata.

\subsection{Performance}

\begin{longtable}[h]{l r r}
  \caption{Custom benchmark results}\label{tab:CustomBenchmark} \\
  
  \toprule
  benchmark name           & runtime (in s) & with mitigation (in s) \\
  \midrule
  \endhead{}
  single-block-used        & 5.8379         & 5.8170                 \\
  \midrule
  many-blocks-used         & 7.0976         & 9.8458                 \\
  \midrule
  random-blocks-order      & 3.8206         & 4.3806                 \\
  \midrule
  random-half-blocks-order & 3.7602         & 4.5997                 \\
  \midrule
  many-many-blocks-used    & 6.4697         & 9.4870                 \\
  \bottomrule
\end{longtable}

In order to benchmark the changes of the mitigation, we used some custom benchmarks to
test for impact with very targeted programs, results of which are show in
\autoref{tab:CustomBenchmark}, as well as a benchmarking suite for malloc implementations,
mimalloc-bench, which is usually used to compare many different allocators between each
other.~\cite{mimalloc-bench} We slightly modified mimalloc-bench to include our patched
Scudo version, and ran the benchmarks between the latest Scudo version and our patched
version.

While in the worst case (many secondary allocations with many secondary chunks being
allocated at the same time) the performance impact is quite high, in usual usage it should
not matter much. This seems to also be confirmed by the benchmark tests of mimalloc-bench
we ran, some of which are shown in \autoref{tab:MimallocBench}. For the full results check
\autoref{tab:FullMimallocBench} in the Appendix.

\begin{longtable}[h]{l l r r r r r r}
  \caption{Excerpt of mimalloc-bench results}\label{tab:MimallocBench} \\

  \toprule
  test & alloc & time & rss & user & sys & page-faults & -reclaims \\
  \midrule
  \endhead{}
  cfrac & sys & 8.49 & 3171 & 8.49 & 0.0 & 0 & 434 \\
  \midrule
  cfrac & scudo\_fixed & 10.03 & 4645 & 10.03 & 0.0 & 0 & 612 \\
  \midrule
  cfrac & scudo & 9.96 & 4672 & 9.96 & 0.0 & 0 & 611 \\
  \midrule
  espresso & sys & 6.36 & 2540 & 6.35 & 0.01 & 1 & 473 \\
  \midrule
  espresso & scudo\_fixed & 7.05 & 4852 & 7.03 & 0.01 & 0 & 659 \\
  \midrule
  espresso & scudo & 7.02 & 4860 & 7.0 & 0.02 & 0 & 654 \\
  \midrule
  barnes & sys & 3.57 & 58405 & 3.54 & 0.03 & 0 & 16570 \\
  \midrule
  barnes & scudo\_fixed & 3.42 & 59704 & 3.39 & 0.03 & 0 & 16645 \\
  \midrule
  barnes & scudo & 3.77 & 59739 & 3.75 & 0.02 & 0 & 16645 \\
  \bottomrule
\end{longtable}


%%%%%%%%%%%%%%%%%%%%
\chapter{Conclusion}
%%%%%%%%%%%%%%%%%%%%

In this project, we used our previous work and acquired knowledge on the Scudo allocator
and the gef plugin we built for it, to try to find some exploits that work on the Scudo
allocator, similar to the glibc exploitation houses. The source code of the different
exploitation scripts, the used exploitable programs, as well as some utilities to
facilitate exploit development are available at
\url{https://github.com/SIELA1915/house-of-scudo}.

We developed a python library, that takes advantage of the speed of native C code to take
on the task of brute-forcing the cookie, given one header leak. Then we found and
implemented two exploits, which are based on the common technique of forging headers using
the brute-forced cookie.
The Safe Unlink Secondary exploit was coincidentally fixed by a change on March 6th 2023,
with commit $a6e3bb9$. Since the CommitBase exploit however was not fixed, we implemented
a possible mitigation, that would fix both of our exploits. We then evaluated our
mitigation, by checking it indeed fixes both of our exploits and by benchmarking the
performance impact of our mitigation. Since the performance impact was negligible in most
common scenarios, we submitted the mitigation to the official llvm repository in the form
of a pull request on December 13th 2023.


\cleardoublepage{}
\phantomsection{}
\addcontentsline{toc}{chapter}{Bibliography}
\printbibliography{}

% Appendices are optional
\appendix
%%%%%%%%%%%%%%%%%%%%%%%%%%%%%%%%%%%%%%
\chapter{Full mimalloc-bench results}
%%%%%%%%%%%%%%%%%%%%%%%%%%%%%%%%%%%%%%

\begin{longtable}[h]{l l r r r r r r}
  \caption{Full mimalloc-bench results}\label{tab:FullMimallocBench} \\

  \toprule
  test & alloc & time & rss & user & sys & page-faults & -reclaims \\
  \midrule
  \endhead{}
  cfrac & sys & 8.49 & 3171 & 8.49 & 0.0 & 0 & 434 \\
  \midrule
  cfrac & scudo\_fixed & 10.03 & 4645 & 10.03 & 0.0 & 0 & 612 \\
  \midrule
  cfrac & scudo & 9.96 & 4672 & 9.96 & 0.0 & 0 & 611 \\
  \midrule
  espresso & sys & 6.36 & 2540 & 6.35 & 0.01 & 1 & 473 \\
  \midrule
  espresso & scudo\_fixed & 7.05 & 4852 & 7.03 & 0.01 & 0 & 659 \\
  \midrule
  espresso & scudo & 7.02 & 4860 & 7.0 & 0.02 & 0 & 654 \\
  \midrule
  barnes & sys & 3.57 & 58405 & 3.54 & 0.03 & 0 & 16570 \\
  \midrule
  barnes & scudo\_fixed & 3.42 & 59704 & 3.39 & 0.03 & 0 & 16645 \\
  \midrule
  barnes & scudo & 3.77 & 59739 & 3.75 & 0.02 & 0 & 16645 \\
  \midrule
  redis & sys & 9.26 & 7929 & 0.42 & 0.04 & 3 & 1240 \\
  \midrule
  redis & scudo\_fixed & 9.81 & 9587 & 0.44 & 0.05 & 0 & 1482 \\
  \midrule
  redis & scudo & 8.85 & 9648 & 0.4 & 0.06 & 0 & 1480 \\
  \midrule
  larsonN-sized & sys & 3.0 & 374957 & 257.39 & 4.62 & 0 & 106747 \\
  \midrule
  larsonN-sized & scudo\_fixed & 128.05 & 168812 & 94.15 & 152.25 & 15 & 82785 \\
  \midrule
  larsonN-sized & scudo & 129.12 & 168765 & 94.66 & 150.81 & 20 & 82673 \\
  \midrule
  mstressN & sys & 5.05 & 1526414 & 37.09 & 20.58 & 0 & 3288327 \\
  \midrule
  mstressN & scudo\_fixed & 13.81 & 803629 & 119.19 & 108.42 & 0 & 6181076 \\
  \midrule
  mstressN & scudo & 13.7 & 800176 & 115.88 & 105.06 & 0 & 6171572 \\
  \midrule
  rptestN & sys & 4.64 & 162499 & 28.06 & 10.43 & 2 & 87482 \\
  \midrule
  rptestN & scudo\_fixed & 22.79 & 143832 & 60.92 & 40.67 & 3 & 528002 \\
  \midrule
  rptestN & scudo & 21.94 & 145304 & 58.56 & 39.31 & 0 & 526599 \\
  \midrule
  lua & sys & 8.45 & 61516 & 7.81 & 0.57 & 179 & 144312 \\
  \midrule
  lua & scudo\_fixed & 8.44 & 71195 & 7.76 & 0.68 & 2 & 179348 \\
  \midrule
  lua & scudo & 8.52 & 71162 & 7.8 & 0.71 & 0 & 179306 \\
  \midrule
  alloc-test1 & sys & 5.44 & 13859 & 5.43 & 0.01 & 0 & 9342 \\
  \midrule
  alloc-test1 & scudo\_fixed & 5.87 & 15688 & 5.85 & 0.01 & 0 & 11476 \\
  \midrule
  alloc-test1 & scudo & 5.73 & 15740 & 5.72 & 0.01 & 0 & 10898 \\
  \midrule
  alloc-testN & sys & 6.12 & 17233 & 87.82 & 0.04 & 7 & 11993 \\
  \midrule
  alloc-testN & scudo\_fixed & 10.77 & 17896 & 168.24 & 0.06 & 7 & 16705 \\
  \midrule
  alloc-testN & scudo & 12.23 & 17858 & 192.0 & 0.06 & 6 & 19608 \\
  \midrule
  sh6benchN & sys & 0.76 & 335016 & 14.53 & 1.21 & 0 & 83493 \\
  \midrule
  sh6benchN & scudo\_fixed & 39.32 & 371601 & 448.45 & 175.71 & 17 & 137682 \\
  \midrule
  sh6benchN & scudo & 38.37 & 371660 & 450.64 & 179.58 & 15 & 137678 \\
  \midrule
  sh8benchN & sys & 1.47 & 424444 & 47.6 & 0.45 & 0 & 110235 \\
  \midrule
  sh8benchN & scudo\_fixed & 94.2 & 277936 & 749.79 & 4333.51 & 17 & 220871 \\
  \midrule
  sh8benchN & scudo & 95.0 & 278531 & 768.26 & 4356.15 & 25 & 223608 \\
  \midrule
  xmalloc-testN & sys & 1.02 & 126016 & 256.71 & 1.76 & 11 & 109729 \\
  \midrule
  xmalloc-testN & scudo\_fixed & 85.24 & 88415 & 23.92 & 191.14 & 0 & 45575 \\
  \midrule
  xmalloc-testN & scudo & 82.89 & 88232 & 23.99 & 196.19 & 0 & 50801 \\
  \midrule
  cache-scratch1 & sys & 1.82 & 3692 & 1.82 & 0.0 & 0 & 239 \\
  \midrule
  cache-scratch1 & scudo\_fixed & 1.94 & 3909 & 1.94 & 0.0 & 0 & 276 \\
  \midrule
  cache-scratch1 & scudo & 1.88 & 3871 & 1.87 & 0.0 & 0 & 273 \\
  \midrule
  cache-scratchN & sys & 0.09 & 4001 & 4.49 & 0.0 & 0 & 392 \\
  \midrule
  cache-scratchN & scudo\_fixed & 0.1 & 4606 & 4.64 & 0.0 & 0 & 527 \\
  \midrule
  cache-scratchN & scudo & 0.1 & 4685 & 4.74 & 0.0 & 0 & 530 \\
  \midrule
  glibc-simple & sys & 5.43 & 1977 & 5.43 & 0.0 & 0 & 216 \\
  \midrule
  glibc-simple & scudo\_fixed & 7.14 & 3623 & 7.13 & 0.0 & 0 & 365 \\
  \midrule
  glibc-simple & scudo & 7.5 & 3696 & 7.5 & 0.0 & 0 & 380 \\
  \midrule
  glibc-thread & sys & 0.68 & 12809 & 109.56 & 0.03 & 70 & 4955 \\
  \midrule
  glibc-thread & scudo\_fixed & 2.66 & 15507 & 109.55 & 0.04 & 9 & 4254 \\
  \midrule
  glibc-thread & scudo & 2.9 & 15431 & 109.55 & 0.05 & 15 & 4257 \\
  \bottomrule
\end{longtable}


\end{document}
